% !PRINTED 325
\weberpage{325}
%<<CONT
und folglich ist die Formel (12) erwiesen, wenn sie für irgend zwei Primdiskriminanten gilt.

Es ist aber nach (10), wenn $p$ und $p'$ ungerade Primzahlen sind:
\[
(\pm p, \pm p') = \pm\left(\frac{p'}{p}\right), \quad (\pm p', \pm p) = \pm\left(\frac{p'}{p}\right), % [sic] The source prints p'/p in both members; the following prose requires p/p' in the second.
\]
worin rechts das negative Zeichen nur dann steht, wenn links \gesperrt{beide Zeichen negativ}, also $p$ sowohl als $p' \equiv 3 \pmod{4}$ sind. In diesem und nur in diesem Falle ist aber $\left(\dfrac{p'}{p}\right) = -\left(\dfrac{p}{p'}\right)$, sonst $\left(\dfrac{p'}{p}\right) = \left(\dfrac{p}{p'}\right)$. Also ergibt sich für diesen Fall die Formel (13) als richtig.

Ferner ist nach (11) und (3), (5)
\[
(-4, \pm p) = +1, \quad (\pm p, -4) = \pm 1,
\]
\[
(p, \pm 8) = (p, \pm 2) = (p, 2) = (-1)^{\frac{p^2-1}{8}},
\]
\[
(\pm 8, p) = \left(\frac{\pm 8}{p}\right) = (-1)^{\frac{p^2-1}{8}},
\]
\[
[p \equiv 1 \pmod{4}],
\]
\[
(-p, \pm 8) = (-p, \pm 2) = \pm(-p, 2) = -1,
\]
\[
(\pm 8, -p) = \pm\left(\frac{\pm 8}{p}\right) = -1,
\]
\[
[p \equiv 3 \pmod{8}],
\]
woraus für diese Fälle die Formel (13) folgt, die damit allgemein erwiesen ist.

Aus (13) folgt weiter, da jede Quadratzahl eine positive Diskriminante ist, falls $m$ zu $D$ \gesperrt{relativ prim} ist,
\weq{14}{(D, m^2) = 1}
und folglich nach (8):
\weq{15}{(D, m^2 n) = (D, n),}
wenn $m$ und $D$ keinen gemeinschaftlichen Teiler haben.

Wir haben ferner den Satz:
\weq{16}{(D, n) = (D, n'),}
wenn $n \equiv n' \pmod{D}$.

Der Satz ist richtig, wenn $n$ und $D$ einen gemeinschaftlichen Teiler haben, und nach (12) ist er allgemein erwiesen, wenn er für jede Primdiskriminante $D$ gilt. Für diese ergibt er sich aber sofort aus (10), (11).
